\documentclass[12pt,letterpaper]{article}
\usepackage[utf8]{inputenc}
\usepackage[letterpaper, margin=1in, bottom=2.5in, top=2.2in]{geometry}
\usepackage{fancyhdr}
\setlength{\headheight}{15pt}
\setlength{\footskip}{72pt}

\usepackage{xcolor}
\usepackage{amsmath}
\usepackage{anyfontsize}
\usepackage{graphicx}
\usepackage{tikz}
\usepackage{listings}
\usepackage{sourcesanspro}
\renewcommand{\familydefault}{\sfdefault}
\usepackage{fontawesome5}
\usepackage{titlesec}
\usepackage{setspace}
\usepackage{hyperref}
\usepackage{mdframed}
\usepackage{qrcode}
\usetikzlibrary{calc,shapes,positioning}
\usepackage{eso-pic}

% Command for inline code highlighting
\newcommand{\code}[1]{\texttt{\textcolor{accentColor}{#1}}}

\AtBeginDocument{\color{primaryColor}}

% Colors for Python theme
\definecolor{bgColor}{RGB}{13, 17, 23}  % Dark background
\definecolor{primaryColor}{RGB}{255, 255, 255}  % White text
\definecolor{accentColor}{RGB}{68, 171, 255}  % Python blue for main accent
\definecolor{accentColor2}{RGB}{200, 160, 35}  % Mustard yellow for highlighting concepts
\definecolor{pythonSecondary}{RGB}{180, 180, 180}  % Light gray
\definecolor{secondaryColor}{RGB}{230, 235, 240}  % Almost white gray
\definecolor{terminalBg}{RGB}{22, 22, 22}  % Terminal background
\definecolor{terminalFrame}{RGB}{40, 40, 40}  % Terminal frame
\definecolor{lineNumberColor}{RGB}{100, 100, 100}  % Line number color
\definecolor{dividerColor}{RGB}{120, 130, 150}  % Divider color

% Code colors
\definecolor{codeTextColor}{RGB}{255, 255, 255}  % Basic text
\definecolor{codeKeywordColor}{RGB}{68, 171, 255}  % Keywords in Python blue
\definecolor{codeCommentColor}{RGB}{98, 150, 85}  % Comments 
\definecolor{codeStringColor}{RGB}{200, 160, 35}  % Strings in mustard yellow
\definecolor{codeClassColor}{RGB}{78, 201, 176}  % Classes
\definecolor{codeFunctionColor}{RGB}{200, 160, 35}  % Functions in mustard yellow
\definecolor{codeNumberColor}{RGB}{181, 206, 168}  % Numbers

\pagecolor{bgColor}

\hypersetup{
    colorlinks=true,
    linkcolor=accentColor,
    filecolor=accentColor,
    urlcolor=accentColor,
}

% Define Python language for listings
\lstdefinelanguage{Python}{
  keywords={and, as, assert, break, class, continue, def, del, elif, else, except, finally, for, from, global, if, import, in, is, lambda, nonlocal, not, or, pass, raise, return, try, while, with, yield, async, await},
  sensitive=true,
  comment=[l]{\#},
  morestring=[b]',
  morestring=[b]"
}

% Listings configuration
\lstset{
  language=Python,
  basicstyle=\ttfamily\bfseries\color{codeTextColor},
  backgroundcolor=\color{terminalBg},
  commentstyle=\color{codeCommentColor},
  keywordstyle=\color{codeKeywordColor},
  stringstyle=\color{codeStringColor},
  numberstyle=\color{lineNumberColor},
  breaklines=true,
  breakatwhitespace=true,
  tabsize=4,
  showstringspaces=false,
  frame=none,
  xleftmargin=15pt,
  xrightmargin=0pt,
  aboveskip=5pt,
  belowskip=5pt,
  numbers=left,
  numbersep=8pt,
  extendedchars=true,
  keepspaces=true,
  columns=flexible,
  lineskip=2pt,
  emph={[2]print,len,type,id,dict,list,tuple,set,str,int,float,bool,None,True,False},
  emphstyle={[2]\color{codeFunctionColor}}
}

\newenvironment{macterminal}{%
    \begin{mdframed}[
        linecolor=terminalFrame,
        backgroundcolor=terminalBg,
        roundcorner=5pt,
        skipabove=5pt,
        skipbelow=5pt,
        linewidth=1pt,
        innertopmargin=5pt,
        frametitle={%
            \tikz[baseline=(current bounding box.east), outer sep=0pt]{
                \fill[red!80!black] (0,0) circle (5pt);
                \fill[yellow!80!black] (0.7,0) circle (5pt);
                \fill[green!70!black] (1.4,0) circle (5pt);
            }
        },
        frametitlealignment=\raggedright,
        frametitleaboveskip=8pt,
        frametitlebelowskip=0pt,
    ]
}{%
    \end{mdframed}%
}

\newcommand{\verspace}{\vspace{5pt}}
\newcommand{\languagetag}[1]{
    \begin{tikzpicture}[baseline]
        \node[fill=accentColor, text=primaryColor, rounded corners=5pt, inner sep=7pt] {
            {\normalsize\textbf{#1}}
        };
    \end{tikzpicture}
}

% Title formatting
\titleformat{\section}
  {\LARGE\bfseries\color{primaryColor}}
  {\thesection. }
  {0pt}
  {}
  []

\titleformat{\subsection}
  {\Large\bfseries\color{accentColor}}
  {\thesubsection. }
  {0pt}
  {}
  []

\titleformat{\subsubsection}
  {\large\bfseries\color{accentColor2}}
  {\thesubsubsection. }
  {0pt}
  {}
  []

\titlespacing*{\section}{0pt}{18pt}{8pt}
\titlespacing*{\subsection}{0pt}{12pt}{5pt}
\titlespacing*{\subsubsection}{0pt}{8pt}{3pt}

% Header and footer setup
\pagestyle{fancy}
\fancyhf{}

\fancyhead[L]{
    \begin{tikzpicture}[remember picture, overlay]
        \fill[accentColor, rounded corners=3pt] (0,0) rectangle (2.2cm,0.7cm);
        \node[text=primaryColor, font=\bfseries] at (1.1cm,0.35cm) {PYTHON};
        \fill[primaryColor] (2.35cm,0.1cm) rectangle (2.40cm,0.6cm);
        \node[text=primaryColor, font=\bfseries, anchor=west] at (2.35cm,0.35cm) {NAMESPACES};
    \end{tikzpicture}
}

\renewcommand{\headrulewidth}{0pt}
\renewcommand{\footrulewidth}{0.2pt}
\renewcommand{\footruleskip}{1cm}

\fancyfoot[C]{
    \vspace*{0.1cm}
    \noindent
    \begin{minipage}{\textwidth}
        \begin{flushleft}
            % Profile photo
            \raisebox{0.7cm}{
            \begin{tikzpicture}[baseline]
                \path[fill=bgColor] (0,0) circle (0.6cm);
                \clip (0,0) circle (0.6cm);
                \node at (0,0) {
                    \includegraphics[width=1.2cm,height=1.2cm]{../../Images/profile-image.jpeg}
                };
            \end{tikzpicture}
            }
            % Profile information
            \begin{minipage}[b]{0.5\textwidth}
                {\large\bfseries\color{primaryColor}Alejandro Sánchez Yalí}
                \par\vspace{1pt}
                {\small\color{secondaryColor}Software Developer | AI \& Blockchain Enthusiast}
                \par\vspace{1pt}
                {\small\color{accentColor2}\faGlobe\hspace{5pt}\color{secondaryColor}www.asanchezyali.com}
            \end{minipage}
        \end{flushleft}
        \vspace{6pt}
    \end{minipage}
}

\fancypagestyle{plain}{
    \fancyhf{}
    \renewcommand{\headrulewidth}{0pt}
    \renewcommand{\footrulewidth}{0pt}
}

\renewcommand{\labelitemi}{\textcolor{accentColor}{$\bullet$}}
\renewcommand{\labelitemii}{\textcolor{accentColor2}{$\circ$}}

\usepackage{relsize}
\AtBeginDocument{\relsize{1}}

% Command for elegant QR
\newcommand{\elegantqr}[2]{
    \qrcode[height=2.5cm]{#1}
    \\[0.1cm]
    {\hspace{0.2cm}\color{primaryColor}\small #2\par}
}
